% gen_server Startup Sequence Diagram (Standalone)
% Compile with: xelatex gen_server_startup.tex
\documentclass[border=10pt]{standalone}

% ============================================================
% Packages
% ============================================================
\usepackage{fontspec}
\usepackage{xeCJK}
\setCJKmainfont[BoldFont=STHeiti, ItalicFont=STKaiti]{STSong}
\setCJKsansfont{STHeiti}
\setCJKmonofont{STFangsong}

\usepackage{xcolor}
\usepackage{tikz}
\usepackage{pifont}  % for \ding symbols (circled numbers)

\usetikzlibrary{arrows.meta, positioning, shapes.geometric, fit, calc, backgrounds, decorations.pathreplacing}

% ============================================================
% Color Definitions (same as gen_server_architecture.tex)
% ============================================================
\definecolor{erlred}{HTML}{A4070A}
\definecolor{erlblue}{HTML}{2B5EA7}
\definecolor{erlgreen}{HTML}{2E7D32}
\definecolor{erlyellow}{HTML}{F9A825}
\definecolor{erlpurple}{HTML}{6A1B9A}
\definecolor{erlmodule}{HTML}{D84315}  % dedicated color for ?MODULE
\definecolor{erlinit}{HTML}{00897B}    % dedicated vivid teal color for init/1
\definecolor{erlpid}{HTML}{E65100}     % dedicated color for Pid/Name return

% ============================================================
% Startup Sequence Diagram
% ============================================================
\begin{document}

\begin{tikzpicture}[scale=0.68, every node/.style={scale=0.68},
    quadrant/.style={minimum width=6.5cm, minimum height=4.5cm, text width=5.8cm, align=left, font=\small},
    api/.style={font=\scriptsize\ttfamily, inner sep=2pt},
]
    % ---- Outer frames ----
    % Left half — Caller Process
    \node[draw=erlpurple!60, fill=erlpurple!3, rounded corners=6pt,
          minimum width=6.6cm, minimum height=11.0cm] (leftframe) at (-4.8, -0.8) {};
    \node[font=\small\bfseries\color{erlpurple}, anchor=south] at (leftframe.north) {调用者进程(Caller)};
\node[font=\small\ttfamily\bfseries\itshape\color{black}, anchor=north west] at ([xshift=3pt, yshift=-1pt]leftframe.north west) {counter\_server.erl};

    % Right half — New Server Process
    \node[draw=erlblue!60, fill=erlblue!3, rounded corners=6pt,
          minimum width=6.6cm, minimum height=11.0cm] (rightframe) at (6.4, -0.8) {};
    \node[font=\small\bfseries\color{erlblue}, anchor=south] at (rightframe.north) {新服务进程(New Server Process)};
\node[font=\small\ttfamily\bfseries\itshape\color{black}, anchor=north west] at ([xshift=3pt, yshift=-1pt]rightframe.north west) {gen\_server.erl};

    % ---- Left-Top: Client Call Site ----
    \node[draw=erlpurple, fill=erlpurple!8, rounded corners=5pt,
          minimum width=6.2cm, minimum height=3.2cm, text width=5.8cm, align=left] (lt) at (-4.8, 2.5) {};
    \node[font=\small\bfseries\color{erlpurple}, anchor=north west] at ([xshift=3pt, yshift=-3pt]lt.north west) {启动调用};
    % Two forms of start_link
    \node[draw=erlred!70, fill=erlred!8, rounded corners=3pt,
          font=\scriptsize\ttfamily\bfseries, inner sep=3pt, text width=5.4cm, align=left, anchor=north]
        (startcalls) at ([yshift=-18pt]lt.north)
        {\textcolor{erlred!80!black}{%
         \ding{172} 匿名启动（返回Pid）：}\\
         gen\_server:start\_link(\textcolor{erlmodule}{\textbf{?MODULE}}, Args, [])\\[4pt]
         \textcolor{erlred!80!black}{\ding{173} 注册名启动（返回Pid+注册名字）：}\\
         gen\_server:start\_link(\\
         \quad \textcolor{erlpid}{\textbf{\{local, counter\}}}, \textcolor{erlmodule}{\textbf{?MODULE}}, Args, [])
        };

    % ---- Left-Bottom: Return Value ----
    \node[draw=erlpid!80, fill=erlpid!8, rounded corners=5pt,
          minimum width=6.2cm, minimum height=3.2cm, text width=5.8cm, align=left] (lb) at (-4.8, -2.8) {};
    \node[font=\small\bfseries\color{erlpid!80!black}, anchor=north west] at ([xshift=3pt, yshift=-3pt]lb.north west) {启动返回值};
    \node[font=\scriptsize\ttfamily, text width=5.4cm, align=left, anchor=north west] at ([xshift=6pt, yshift=-18pt]lb.north west) {
        \textcolor{erlpid!80!black}{\textbf{成功：}}\\
        \{ok, \textcolor{erlpid}{\textbf{Pid}}\} = start\_link(...)\\
        \textcolor{gray}{\%\% Pid = <0.123.0>}\\[4pt]
        \textcolor{erlpid!80!black}{\textbf{注册名后可直接使用名字：}}\\
        gen\_server:call(\textcolor{erlpid}{\textbf{counter}}, get\_count)\\[4pt]
        \textcolor{erlred!80!black}{\textbf{失败：}}\\
        \{error, Reason\}\\
        ignore
    };

    % ---- Right-Top: Startup Sequence (inside new process) ----
    \node[draw=erlblue, fill=erlblue!10, rounded corners=5pt,
          minimum width=6.2cm, minimum height=3.2cm, text width=5.8cm, align=left] (rt) at (6.4, 2.5) {};
    \node[font=\small\bfseries\color{erlblue}, anchor=north west] at ([xshift=3pt, yshift=-3pt]rt.north west) {启动序列（新进程内）};
    \node[font=\scriptsize, text width=5.4cm, align=left, anchor=north west] at ([xshift=6pt, yshift=-18pt]rt.north west) {
        \textcolor{erlblue!80!black}{\ding{172}} \texttt{proc\_lib:start\_link/5}\\
        \quad $\rightarrow$ 创建新进程 (\textcolor{erlpid}{\textbf{Pid}})\\[3pt]
        \textcolor{erlblue!80!black}{\ding{173}} \texttt{init\_it/6}\\
        \quad $\rightarrow$ 注册名字 \textcolor{erlpid}{\textbf{\{local, Name\}}}\\
        \quad $\rightarrow$ 调用 \textcolor{erlinit}{\textbf{\texttt{Mod:init(Args)}}}\\[3pt]
        \textcolor{erlblue!80!black}{\ding{174}} \texttt{proc\_lib:init\_ack(\textcolor{erlpid}{\textbf{\{ok, Pid\}}})}\\
        \quad $\rightarrow$ 通知调用者启动成功
    };

    % ---- Right-Bottom: Process Registration & State ----
    \node[draw=erlgreen!70!black, fill=erlgreen!5, rounded corners=5pt,
          minimum width=6.2cm, minimum height=3.2cm, text width=5.8cm, align=left] (rb) at (6.4, -2.8) {};
    \node[font=\small\bfseries\color{erlgreen!70!black}, anchor=north west] at ([xshift=3pt, yshift=-3pt]rb.north west) {进程标识（启动完成后）};
    \node[font=\scriptsize, text width=5.4cm, align=left, anchor=north west] at ([xshift=6pt, yshift=-18pt]rb.north west) {
        \textcolor{erlpid!80!black}{\textbf{Pid 标识：}}\\
        \texttt{Pid = <0.123.0>}\\
        \texttt{gen\_server:call(\textcolor{erlpid}{\textbf{Pid}}, Req)}\\[4pt]
        \textcolor{erlpid!80!black}{\textbf{注册名标识：}}\\
        \texttt{register(\textcolor{erlpid}{\textbf{counter}}, Pid)}\\
        \texttt{gen\_server:call(\textcolor{erlpid}{\textbf{counter}}, Req)}\\[4pt]
        \textcolor{gray}{\textbf{全局名标识：}}\\
        \texttt{\{global, Name\} | \{via, Mod, Name\}}
    };

    % ==== Arrows ====

    % Arrow A: start_link call — Caller to New Process (spawn)
    \draw[-{Stealth[length=6pt]}, ultra thick, erlred!80!black, line width=1.6pt, densely dashed]
        ([yshift=6mm]lt.east) -- node[above, font=\scriptsize, align=center, text=erlred!80!black]
        {\textcolor{black}{\ding{172}} \texttt{start\_link(\textcolor{erlmodule}{\textbf{?MODULE}}, Args)}\\[1pt]
         {\tiny 同步调用，阻塞等待 init 完成}} ([yshift=6mm]rt.west);

    % Arrow B: init_ack — New Process back to Caller ({ok, Pid})
    \draw[-{Stealth[length=6pt]}, ultra thick, erlpid!90!black, line width=1.6pt]
        ([yshift=-6mm]rt.west) -- node[below, font=\scriptsize, align=center, text=erlpid!80!black]
        {\textcolor{black}{\ding{173}} 返回 \texttt{\textcolor{erlpid}{\textbf{\{ok, Pid\}}}}\\[1pt]
         {\tiny init/1 返回 \{ok, State\} 后通过 init\_ack 通知}} ([yshift=-6mm]lt.east);

    % Arrow C: Name registration (inside right side)
    \draw[-{Stealth[length=5pt]}, thick, erlpid!70, line width=1.2pt, dashed]
        ([xshift=-5mm]rt.south) -- node[left, font=\scriptsize, text=erlpid!70, align=right]
        {\textcolor{black}{\ding{174}} 注册名字\\{\tiny\texttt{register(Name, Pid)}}} ([xshift=-5mm]rb.north);

    % Arrow D: Caller receives Pid, can now use it
    \draw[-{Stealth[length=5pt]}, thick, erlgreen!70!black, line width=1.2pt]
        ([xshift=0mm]lt.south) -- node[left, font=\scriptsize, text=erlgreen!70!black, align=right]
        {\textcolor{black}{\ding{175}} 获得 Pid/Name\\{\tiny 后续用于 call/cast}} ([xshift=0mm]lb.north);

    % Arrow E: Runtime usage — from return value box to process identity
    \draw[-{Stealth[length=5pt]}, thick, erlblue, line width=1.0pt, dotted]
        (lb.east) -- node[api, above, align=center, text width=3.5cm]
        {\scriptsize \textcolor{black}{\ding{176}} \textcolor{erlblue}{运行时通信}\\[1pt]
         \tiny \texttt{call(Pid, Req)} 或 \texttt{call(Name, Req)}} (rb.west);

    % ---- Dashed separator lines below inner boxes ----
    \draw[dashed, erlpurple!40, line width=0.6pt]
        ([yshift=-6pt]lb.south west) -- ([yshift=-6pt]lb.south east);
    \draw[dashed, erlblue!40, line width=0.6pt]
        ([yshift=-6pt]rb.south west) -- ([yshift=-6pt]rb.south east);

    % Left bottom annotation
    \node[font=\tiny\color{erlred!70!black}, anchor=north west, text width=5.8cm, align=left] at ([xshift=6pt, yshift=-10pt]lb.south west)
        {\ding{72} \texttt{start\_link} 是\textbf{同步}的：调用者阻塞直到 \textcolor{erlinit}{\textbf{init/1}} 完成\\
         \ding{72} 返回的 \textcolor{erlpid}{\textbf{Pid}} 是后续所有通信的句柄};

    % Right bottom annotation
    \node[font=\tiny\color{erlred!70!black}, anchor=north west, text width=5.8cm, align=left] at ([xshift=6pt, yshift=-10pt]rb.south west)
        {\ding{72} 注册名使得无需传递 Pid，直接用原子名通信\\
         \ding{72} \texttt{\{local, Name\}} 仅本节点；\texttt{\{global, Name\}} 跨节点};

\end{tikzpicture}

\end{document}
